% =============================================================================
%  Capítulo 2 — Inteligência Artificial (Parte II)
%  Capítulo único da parte. Seções: o que a IA faz e não faz, como nasce uma
%  pré-análise (BPMN), seis passos com a tela (lista do Assistente, leitura,
%  decisão, Consulta, Memória e Desempenho, tradução), selos e rotas, a
%  tradução por dentro, boas práticas e mensagens.
%
%  As telas de img/ são ESQUEMAS provisórios (rodapé "CAPTURA PENDENTE"): trocar
%  pelas capturas reais, na mesma resolução, e recalibrar os marcadores x/y.
% =============================================================================
\chapter{Inteligência Artificial}\label{cap:ia}

\section{O que a IA faz e o que não faz}

O Hanaro usa inteligência artificial em duas frentes: o \textbf{assistente}, que
prepara pré-análises e responde a perguntas sobre casos já tratados, e a
\textbf{tradução} dos relatórios para o inglês. Em ambas, a IA \emph{propõe} e uma
pessoa \emph{decide}. A \autoref{tab:funcoes} resume as funções.

\begin{table}[H]
\caption{Funções de IA do Hanaro.}\label{tab:funcoes}
\guiazebra
\small
\noindent\begin{tabularx}{\textwidth}{@{}L{2.7cm}X L{3.4cm}@{}}
    \toprule
    \guiacabtabela
    \thead{Função} & \thead{O que faz} & \thead{Quem usa} \\
    \midrule
    \textbf{Pré-análise} & Lê os fatos de um caso, busca casos parecidos na memória e
        propõe os campos da análise e as tarefas. & Analista e Gestor \\
    \textbf{Consulta} & Responde a uma pergunta em texto livre usando só análises já
        aprovadas, sempre com as fontes. & Analista e Gestor, quando ligada \\
    \textbf{Memória} & Reúne as análises aprovadas que servem de exemplo ao
        assistente. & Gestor e Administrador \\
    \textbf{Desempenho} & Mostra o que as pessoas decidiram sobre as propostas. &
        Gestor e Administrador \\
    \textbf{Tradução} & Traduz o texto de um relatório do português para o inglês
        antes do download. & Quem emite o documento \\
\end{tabularx}
\end{table}

\begin{remark}
\textbf{A IA nunca decide sozinha.} Nada do que o assistente propõe vira análise,
relatório ou tarefa até que uma pessoa aceite. O assistente também não altera dados do
Hanaro por conta própria e não chama o Hanaro: é o Hanaro que lhe entrega o caso e
busca a resposta.
\end{remark}

O que o assistente \textbf{não} faz: não afirma uma causa quando não há precedente (só
traz os fatos, o checklist 4M e perguntas para quem analisa); não cita fonte, número ou
data que não estejam nos dados; e não repete como certeza o que é hipótese. Toda causa
proposta é marcada como \emph{hipótese}, para ser confirmada por quem conhece a linha.

\begin{remark}[info]
\textbf{Ativação gradual.} O assistente, a Consulta e a tradução dependem do servidor
de IA e são ligados por ambiente. Quando estão desligados, o Hanaro avisa na própria
tela, as pré-análises já criadas continuam disponíveis e o restante do sistema segue
normal.
\end{remark}

\section{Como nasce uma pré-análise}

A pré-análise começa em uma \textbf{regra de alerta} configurada para pedir ajuda do
assistente. Quando a regra dispara sobre uma ocorrência, o Hanaro entrega o caso ao
agente e aguarda. O agente reúne os fatos, procura casos parecidos na memória e prepara
uma proposta, que passa por uma conferência antes de chegar à pessoa. A
\autoref{fig:ia} mostra o fluxo; os números nos quadros vermelhos levam aos passos
deste capítulo.

\begin{figure}[H]
\centering
\begin{tikzpicture}[x=1cm,y=1cm,
        bpmn tarefa/.append style={text width=4.0cm,minimum height=0.8cm}]
    \useasboundingbox (0,0) rectangle (16,-13.7);
    \bpmnraia{0}{5}{13.7}{Analista ou Gestor}
    \bpmnraia{5}{5.6}{13.7}{Hanaro}
    \bpmnraia{10.6}{5.4}{13.7}{Agente}

    % ---- Hanaro ----
    \node[bpmn inicio] (ini) at (7.8,-1.25) {};
    \node[bpmn rotulo,anchor=west] at ([xshift=4pt]ini.east) {Alerta dispara};
    \node[bpmn tarefa] (h1) at (7.8,-2.4) {Regra de alerta pede pré-análise};
    \node[bpmn tarefa] (h2) at (7.8,-3.7) {Entregar o caso ao agente};
    \node[bpmn tarefa] (h3) at (7.8,-8.0) {Mostrar o selo \emph{Pré-análise pronta}};
    \node[bpmn tarefa] (h4) at (7.8,-11.7) {Criar análise, relatório e tarefas aceitos};

    % ---- agente ----
    \node[bpmn tarefa] (a1) at (13.3,-3.7) {Reunir os fatos e buscar casos parecidos};
    \bpmngateway{g1}{13.3,-5.0}
    \node[bpmn rotulo,anchor=south west,inner sep=0pt] at ([xshift=2pt,yshift=-1pt]g1.north east) {Há precedente?};
    \node[bpmn tarefa,text width=2.3cm,minimum height=1.5cm] (a2) at (11.95,-6.5)
        {Adaptar campos e tarefas dos casos anteriores};
    \node[bpmn tarefa,text width=2.3cm,minimum height=1.5cm] (a3) at (14.65,-6.5)
        {Só fatos, checklist 4M e perguntas};
    \node[bpmn tarefa,text width=4.3cm] (a4) at (13.3,-8.0)
        {Conferir: recusar fonte ou número inventado};
    \node[bpmn tarefa] (a5) at (13.3,-11.7) {Guardar a análise aprovada na memória};
    \node[bpmn fim] (fim) at (13.3,-13.0) {};

    % ---- pessoa ----
    \bpmnpasso{p1}{2.5,-9.4}{1}{Abrir a lista do Assistente}
    \bpmnpasso{p2}{2.5,-10.55}{2}{Ler a pré-análise}
    \bpmnpasso{p3}{2.5,-11.7}{3}{Aceitar, editar ou rejeitar}

    % ---- fluxos ----
    \draw[bpmn fluxo] (ini) -- (h1);
    \draw[bpmn fluxo] (h1) -- (h2);
    \draw[bpmn fluxo] (h2) -- (a1);
    \draw[bpmn fluxo] (a1) -- (g1);
    \draw[bpmn fluxo] (g1.west) -| node[bpmn rotulo,anchor=north east,pos=0.15] {Sim} (a2.north);
    \draw[bpmn fluxo] (g1.east) -| node[bpmn rotulo,anchor=north west,pos=0.15] {Não} (a3.north);
    \draw[bpmn fluxo] (a2.south) -- (a2.south |- a4.north);
    \draw[bpmn fluxo] (a3.south) -- (a3.south |- a4.north);
    \draw[bpmn fluxo] (a4.west) -- (h3.east);
    \draw[bpmn fluxo] (h3.south) |- (p1.east);
    \draw[bpmn fluxo] (p1) -- (p2);
    \draw[bpmn fluxo] (p2) -- (p3);
    \draw[bpmn fluxo] (p3.east) -- (h4.west);
    \draw[bpmn fluxo] (h4.east) -- (a5.west);
    \draw[bpmn fluxo] (a5) -- (fim);
\end{tikzpicture}
\caption{Da regra de alerta à decisão da pessoa, em BPMN. Os números indicam os passos
deste capítulo.}
\label{fig:ia}
\end{figure}

\begin{remark}[info]
\textbf{Para TI.} O gatilho é a opção de pré-análise da regra de alerta (nível
\tech{agent\_preanalysis}). A entrega ao agente é feita pelo worker, com nova tentativa;
o resultado é buscado pelo Hanaro. Se o agente falha, só a pré-análise atrasa.
\end{remark}

\section{As peças do agente}

O agente é um serviço à parte do Hanaro, no servidor de IA. Ele não acessa o banco do Hanaro: só recebe os casos que o Hanaro entrega, responde e guarda a sua própria memória. A \autoref{fig:agente-componentes} mostra as seis peças do agente e como elas se ligam, no formato de um diagrama de componentes da UML. Cada caixa é uma peça; o círculo é o que a peça oferece, o meio círculo é o que ela precisa de outra, e as pequenas caixas na borda são as portas por onde o serviço conversa com o mundo de fora.

\begin{table}[H]
\caption{As peças do agente de IA.}\label{tab:agente-pecas}
\guiazebra
\small
\noindent\begin{tabularx}{\textwidth}{@{}L{3.2cm}X@{}}
    \toprule
    \guiacabtabela
    \thead{Peça} & \thead{O que faz} \\
    \midrule
    \textbf{API HTTP} & É a porta de entrada do serviço. Recebe os pedidos do Hanaro, protegidos por chave de API, para criar jobs, buscar casos e conversar. \\
    \textbf{Fila e worker} & Guarda os pedidos de pré-análise em uma fila e executa os pedidos em ordem. O mesmo pedido enviado duas vezes não vira duas pré-análises. \\
    \textbf{Grafo de pré-análise} & O roteiro da pré-análise: reunir os fatos, buscar casos parecidos, dar nota, escolher a rota, adaptar a proposta, conferir e finalizar. \\
    \textbf{Memória de casos} & Reúne as análises aprovadas e as encontra por significado e por palavras. A reordenação dos resultados é opcional, e cada análise aceita aumenta a memória. \\
    \textbf{Cliente de modelos} & Fala com o servidor de modelos: gera textos e vetores, e exige respostas em formato estruturado quando o roteiro precisa. \\
    \textbf{Chat de consulta} & Responde às perguntas da Consulta, com as fontes, usando só os casos da memória. \\
    \bottomrule
\end{tabularx}
\end{table}

\begin{remark}[info]
\textbf{Para TI.} A fila de jobs fica no PostgreSQL do próprio agente, o mesmo banco que guarda a memória de casos (com pgvector). O agente nunca chama o Hanaro: quem entrega os pedidos e busca os resultados é o worker do Hanaro.
\end{remark}

\begin{figure}[H]
\centering
\includegraphics[width=\dimexpr\textheight-3.2cm\relax,height=\textwidth,keepaspectratio,angle=90]{img/diagramas/componentes-agent-service}
\caption{Peças do agente de IA, em diagrama de componentes.}\label{fig:agente-componentes}
\end{figure}

\section{Rotas, selos e situações}

Cada pré-análise indica a \textbf{rota} que o agente seguiu, e a rota diz o quanto a
proposta pode ser aproveitada. A \autoref{tab:rotas} traz as três rotas; a
\autoref{tab:situacoes}, as situações em que uma pré-análise pode estar.

\begin{table}[H]
\caption{Rotas da pré-análise.}\label{tab:rotas}
\guiazebra
\small
\noindent\begin{tabularx}{\textwidth}{@{}L{3.2cm}X@{}}
    \toprule
    \guiacabtabela
    \thead{Rota} & \thead{O que significa} \\
    \midrule
    \textbf{Recorrência} & O item já apareceu em casos anteriores do Hanaro. A proposta
        parte do que foi feito antes, e o relatório pode se ligar ao plano já
        existente. \\
    \textbf{Caso parecido} & Há casos parecidos na memória. A proposta se apoia em até
        três deles, sem tratá-los como a mesma causa. \\
    \textbf{Sem precedente} & Não há nada parecido na memória. O assistente traz só os
        fatos, o checklist 4M e perguntas para quem analisa; qualquer causa mostrada é
        uma hipótese a confirmar. \\
\end{tabularx}
\end{table}

\begin{table}[H]
\caption{Situações de uma pré-análise.}\label{tab:situacoes}
\guiazebra
\small
\noindent\begin{tabularx}{\textwidth}{@{}L{3.9cm}X@{}}
    \toprule
    \guiacabtabela
    \thead{Situação} & \thead{O que significa} \\
    \midrule
    \textbf{Aguardando decisão} & Está pronta e ninguém decidiu ainda. É a que conta no
        menu e na lista. \\
    \textbf{Aplicada} & As decisões foram aplicadas: a análise, o relatório e as
        tarefas aceitas foram criados. \\
    \textbf{Aplicada em parte} & Parte das decisões foi aplicada; o restante não pôde
        ser. \\
    \textbf{Dispensada} & Alguém a descartou, com motivo. \\
    \textbf{Substituída} & O agente preparou uma versão mais nova, e esta ficou como
        histórico. \\
\end{tabularx}
\end{table}

O selo \textbf{Pré-análise pronta} aparece na tela \textbf{Alertas}, ao lado do alerta
que gerou a pré-análise, e leva direto a ela. Enquanto o agente trabalha, a lista mostra
\emph{Preparando…}; se houver versão nova, o aviso \emph{Nova versão em preparo}.

% ---------------------------------------------------------------------------
%  Os seis passos. Coordenadas dos marcadores: esquemas provisórios.
% ---------------------------------------------------------------------------
\passo[0.012/0.140, 0.182/0.090, 0.182/0.190, 0.182/0.360]{01-assistente-lista}{Abrir a lista do Assistente}{%
    A lista reúne as pré-análises preparadas para o seu trabalho. O contador ao lado do
    menu mostra quantas estão prontas e aguardando você, e é atualizado a cada minuto.
    \begin{elementos}
        \item \caminho{Assistente}, no menu lateral: abre a lista de pré-análises.
        \item \textbf{Seções}: \botao{Pré-análises}, \botao{Consulta},
              \botao{Memória} e \botao{Desempenho}. Aparecem só as que o seu perfil
              pode usar.
        \item \textbf{Abas e filtros}: \botao{Aguardando decisão},
              \botao{Decididas} e \botao{Todas}; filtros por gatilho, rota, linha e
              período.
        \item \textbf{Pré-análise}: rota, item, linha, valor e o gatilho que a
              disparou. Um clique abre a pré-análise (passo~2).
    \end{elementos}
}

\passo[0.182/0.120, 0.182/0.400, 0.182/0.660, 0.592/0.660]{02-preanalise-leitura}{Ler a pré-análise}{%
    Leia a pré-análise de cima para baixo antes de decidir. Ela mostra o que o
    assistente viu e de onde tirou cada sugestão.
    \begin{elementos}
        \item \textbf{Rota e \emph{O que aconteceu}}: a faixa da rota explica em que a
              proposta se apoia. Os fatos trazem item, linha, ocorrências, período,
              valor e a comparação com a mediana das janelas anteriores.
        \item \textbf{Proposta para a análise}: título, descrição, tipo de defeito,
              família da causa (4M), causa (hipótese) e contenção imediata.
        \item \textbf{Casos parecidos}: análises aprovadas de outros casos, com os
              motivos da semelhança. \botao{Não é parecido} faz o assistente dar menos
              peso ao caso nas próximas buscas.
        \item \textbf{Por que o assistente propõe isso} e \textbf{Fontes}: a
              justificativa de cada campo e a origem do dado (caso, relatório, regra ou
              dado do alerta).
    \end{elementos}
}

\begin{remark}
Parecido não é a mesma causa. Confira os casos citados e o histórico da linha antes de
aproveitar uma sugestão. Na rota \emph{Sem precedente}, a causa é só uma hipótese.
\end{remark}

\passo[0.182/0.120, 0.182/0.360, 0.182/0.600, 0.182/0.800]{03-preanalise-decisao}{Aceitar, editar ou rejeitar}{%
    Nenhuma sugestão é aplicada até você decidir sobre ela.
    \begin{elementos}
        \item \textbf{Cada campo}: \botao{Aceitar}, \botao{Editar} ou
              \botao{Rejeitar}. \botao{Desfazer} volta atrás e \botao{Aceitar tudo}
              aceita os campos de uma vez. Ao rejeitar, o campo \emph{Por quê?}
              (opcional) registra o motivo: \emph{Está errado}, \emph{Não se aplica},
              \emph{Repete outra coisa}, \emph{Já foi feito} ou \emph{Outro motivo}.
        \item \textbf{Tarefas propostas}: cada tarefa também é aceita, editada ou
              rejeitada. Tarefas só existem dentro de um relatório.
        \item \campo{Criar o relatório e finalizar a análise}: sem marcar esta opção, a
              análise fica como rascunho para você completar. Com ela, informe o título
              e o prazo do relatório. O responsável em branco fica com você.
        \item \botao{Aplicar decisões}: grava o que foi aceito ou editado.
              \botao{Dispensar} descarta a pré-análise, com um motivo.
    \end{elementos}
    Ao final, a pergunta \emph{A pré-análise ajudou?} (\emph{Sim}, \emph{Em parte},
    \emph{Não}) alimenta o Desempenho.
}

\begin{remark}[info]
Se a pré-análise cobre um grupo de ocorrências, marque \emph{Criar o relatório}: o
assistente cria uma análise para cada ocorrência que ainda não tem uma e um relatório
com todas, usando as mesmas decisões. Ocorrências que já tinham análise são mantidas.
\end{remark}

\begin{remark}
Se a tela avisar que \emph{há uma versão mais nova desta pré-análise} ou que ela
\emph{mudou enquanto você a revisava}, clique em \botao{Carregar a versão nova}. As
decisões feitas sobre a versão antiga não são aplicadas sozinhas.
\end{remark}

\passo[0.712/0.090, 0.182/0.780, 0.182/0.220, 0.182/0.490]{04-consulta}{Perguntar na Consulta}{%
    A Consulta responde a perguntas sobre como casos parecidos foram tratados. As
    respostas vêm só de análises já aprovadas e sempre mostram de onde saíram. Está
    disponível a Analistas e Gestores, quando a função está ligada.
    \begin{elementos}
        \item \botao{Histórico} e \botao{Nova conversa}: as conversas ficam salvas só
              para você. \botao{Apagar conversa} as remove de vez.
        \item \textbf{Pergunta}: escreva em texto livre, por exemplo \emph{Como
              resolvemos rebarba no chassi da linha G12?}, e clique em
              \botao{Perguntar}.
        \item \textbf{Resposta}: o texto do assistente, montado a partir dos casos
              encontrados.
        \item \textbf{Fontes}: cada análise citada, com linha, item e data.
              \botao{Abrir análise} leva até ela. \botao{Útil} e
              \botao{Não útil} avaliam a resposta.
    \end{elementos}
}

\begin{remark}[dica]
Sem resultado, tente outras palavras, o código do item ou a linha. A mensagem \emph{Nenhuma
análise parecida no histórico} indica que a memória não tem o que responder, e o
assistente não preenche essa lacuna.
\end{remark}

\passo[0.182/0.120, 0.182/0.290, 0.182/0.400, 0.182/0.800]{05-memoria}{Cuidar da memória e do desempenho}{%
    A \textbf{Memória} guarda as análises aprovadas que o assistente usa como exemplo.
    Os Gestores e os Administradores a mantêm.
    \begin{elementos}
        \item \textbf{Indicadores}: análises aprovadas, na memória, aguardando envio,
              retiradas, na fila de envio e envios com falha.
        \item \textbf{Busca e situação}: procure por título ou item e filtre por
              \emph{Na memória}, \emph{Aguardando envio} ou \emph{Retirada}.
        \item \textbf{Lista}: item, linha e data de cada análise.
              \botao{Retirar da memória} (com motivo opcional) tira a análise dos
              exemplos; \botao{Colocar de volta} a devolve.
        \item \botao{Reenviar tudo}: marca todas as análises para reenvio, em lotes e em
              segundo plano. Use depois de trocar o modelo de embedding ou se a
              memória se perder.
    \end{elementos}
}

\begin{remark}[info]
O \textbf{Desempenho} (\caminho{Assistente\seta Desempenho}) mostra, para um período, o
número de pré-análises decididas, a parcela das propostas aceitas como estavam, o tempo
mediano até a pré-análise ficar pronta, os campos mais rejeitados e seus motivos, a
avaliação de quem usou e o modelo em uso. Os números vêm do que as pessoas decidiram e
não dependem de o assistente estar ligado.
\end{remark}

\passo[0.212/0.140, 0.212/0.370, 0.212/0.520, 0.212/0.670, 0.562/0.840]{06-traducao}{Traduzir o relatório para o inglês}{%
    A tradução acontece ao emitir o documento de um relatório, em \botao{Baixar
    documento}.
    \begin{elementos}
        \item \campo{Idioma do documento}: escolha \emph{English}.
        \item \campo{Traduzir o texto para inglês}: opção com a etiqueta \emph{IA},
              que traduz o texto do documento antes do download. Só aparece quando a
              tradução está ligada.
        \item \textbf{Andamento}: \emph{Aguardando o serviço de tradução}, \emph{Na fila
              de tradução} (com o número de pedidos na sua frente), \emph{Carregando o
              modelo de tradução} e \emph{Traduzindo o texto, X de Y}.
        \item \textbf{Conferência}: se algum texto não passou na conferência, a tela
              lista os que ficarão em português. \botao{Emitir mesmo assim} segue com
              eles.
        \item \botao{Baixar PDF} (ou \botao{Baixar apresentação}): gera o documento.
    \end{elementos}
}

\begin{remark}[dica]
Se a tradução falhar, a tela oferece desligá-la e emitir em português, ou tentar de
novo. O documento em português continua disponível.
\end{remark}

\section{Como a tradução protege o seu relatório}

A tradução não entrega o relatório inteiro ao modelo. O Hanaro \textbf{mascara} antes
o que não pode mudar (números, códigos e marcadores), envia só o texto livre ao serviço
de tradução e, na volta, confere o resultado. Um campo só é gravado em inglês se passar
nas conferências de formato, marcadores e números. O campo que não passa fica em
português, e a tela avisa.

\begin{itemize}
    \item \textbf{Um pedido por vez.} Só um trabalho de tradução fala com o serviço de
          cada vez; os demais esperam na fila, do mais antigo para o mais novo.
    \item \textbf{Falha do serviço.} O Hanaro repete as falhas temporárias de rede. Se
          o serviço seguir indisponível, o trabalho termina como falho e nada
          incompleto é gravado.
    \item \textbf{Revisão humana preservada.} Ao trocar o modelo, as traduções
          automáticas antigas são refeitas, e as revisadas por pessoas permanecem.
\end{itemize}

\begin{remark}[info]
\textbf{Para TI.} A tradução usa o \tech{translation-service}, que fala com um servidor
de modelos de API compatível com a da OpenAI. Ativar, escolher o motor
(\tech{TRANSLATION\_ENGINE}) e trocar o modelo (incrementar
\tech{TRANSLATION\_ENGINE\_REVISION}).
\end{remark}

\section{Propostas de tarefas em um relatório}

Dentro de um relatório, o botão \botao{Propor tarefas} pede ao assistente sugestões de
tarefas, a partir do relatório e dos casos parecidos. Cada proposta mostra o motivo e as
fontes, e traz perguntas para quem analisa. Escolha as que fazem sentido e clique em
\botao{Criar as tarefas escolhidas}; \botao{Dispensar} descarta as propostas. Nada vira
tarefa até que você aceite.

\section{Boas práticas e limites}

\begin{itemize}
    \item \textbf{Leia antes de aceitar.} A proposta é um ponto de partida. Quem conhece
          a linha confirma a causa e a ação.
    \item \textbf{Rejeite com motivo.} O motivo informado alimenta o Desempenho e ajuda
          a equipe a ajustar o assistente.
    \item \textbf{Marque o que não é parecido.} \botao{Não é parecido} melhora as
          próximas buscas.
    \item \textbf{Cuide da memória.} Análises erradas ou pouco confiáveis podem ser
          retiradas dos exemplos; a memória só cresce com análises aprovadas.
    \item \textbf{Sem dado, sem resposta.} Se a memória não tem caso parecido, a
          Consulta diz isso e a pré-análise não inventa uma causa.
    \item \textbf{Modelo pequeno, conferência rígida.} O assistente usa modelos de
          linguagem de porte reduzido, e por isso tudo o que ele propõe passa por
          conferências: fonte inexistente, número ou data fora dos dados, família de
          causa ou tipo de defeito fora dos cadastros e texto acima dos limites do
          Hanaro são recusados.
\end{itemize}

\section{Mensagens e solução de problemas}

\begin{table}[H]
\caption{Mensagens do assistente e da tradução.}\label{tab:mensagens}
\guiazebra
\small
\noindent\begin{tabularx}{\textwidth}{@{}L{5.4cm}X@{}}
    \toprule
    \guiacabtabela
    \thead{Mensagem} & \thead{O que fazer} \\
    \midrule
    \emph{O assistente está desligado neste ambiente} & Siga o trabalho sem ele; as
        pré-análises já criadas continuam disponíveis. Peça à TI para ligá-lo. \\
    \emph{O assistente ainda está preparando esta pré-análise} & Aguarde e atualize a
        tela. \\
    \emph{O assistente não conseguiu preparar esta pré-análise} & Faça a análise normalmente
        na Base de Scrap; avise a TI se se repetir. \\
    \emph{Há uma versão mais nova desta pré-análise} & Clique em \botao{Carregar a
        versão nova} e decida de novo. \\
    \emph{Já existe uma análise desta ocorrência} & Abra a análise existente na Base de
        Scrap. \\
    \emph{O assistente não respondeu. Tente de novo em instantes} & A pergunta continua no
        campo da Consulta; envie de novo. \\
    \emph{O assistente está respondendo outras perguntas} & Aguarde alguns instantes e
        tente de novo. \\
    \emph{Nenhuma análise parecida no histórico} & Reescreva a pergunta com outras
        palavras, o código do item ou a linha. \\
    \emph{Não foi possível traduzir agora} & Desligue a tradução para emitir em
        português, ou tente de novo. \\
    \emph{N texto(s) ficarão em português} & Revise esses trechos ou use \botao{Emitir
        mesmo assim}. \\
\end{tabularx}
\end{table}

\fimdecapitulo
